\documentclass[11pt]{article}

\usepackage{fullpage}
\usepackage{amsmath, amssymb, bm, cite, epsfig, psfrag}
\usepackage{graphicx}
\usepackage{float}
\usepackage{amsthm}
\usepackage{amsfonts}
\usepackage{cite}
\usepackage{hyperref}
\usepackage{pgf,tikz}
\usepackage{enumitem}
\usepackage{mathtools}
\usepackage{siunitx}
\usepackage{../../styles/course}


\begin{document}

\title{Introduction to Hardware Design\\
Loop optimization:  Figures and Code Snippets}
\author{Profs. Sundeep Rangan, Siddharth Garg}
\date{}

\maketitle
\subsection{Quadratic function implementation}

First consider quadratic function:

\begin{systemverilogcode}
always_ff @(posedge clk) begin
  x_reg <= x;
end
always_comb begin
  xsq = x_reg * x_reg;
  xsq_w2 = w2 * xsq;
  x_w1 = w1 * x_reg;
  y = w0 + xsq_w2 + x_w1;
end
\end{systemverilogcode}

 
The block diagram is as shown in Figure~\ref{fig:quad_func}.

\begin{figure}
\centering
\begin{tikzpicture}
    % Input register
    \node (x) {\texttt{x}};
    \node[shape=register, draw, minimum width=0.7cm, minimum height=1cm,
        right=1cm of x,fill=blue!10] (xreg) {};
    \draw[->] (x) -- (xreg.input);
    \draw[-] (xreg.output) -- ++(1.3,0) node [midway,above] {\texttt{x\_reg}}
      coordinate (xreg_split);
    \fill (xreg_split) circle (1.5pt);

    % Multiplier for x^2
    \node [circop, right=0.8cm of xreg_split,fill=green!20] (mul1) {$\times$};
    \draw[->] (xreg_split) -- (mul1);
    \draw[->] (xreg_split) -- ++(0,0.5) -| (mul1);
    \node[below= 0.1cm of mul1] (m1) {$M_1$};
    
    % Multiplier for w2*x^2
    \node [circop, right=1cm of mul1,fill=green!20] (mul2) {$\times$};
    \node [above of=mul2, yshift=0.5cm] (w2) {\texttt{w2}};
    \draw[->] (mul1) -- (mul2);
    \draw[->] (w2) -- (mul2);
    \node[below= 0.1cm of mul2] (m2) {$M_2$};

    % Multiplier for w1*x
    \node [circop, below=1cm of mul2,fill=green!20] (mul3) {$\times$};
    \node [below=1cm of mul3] (w1) {\texttt{w1}};
    \draw[->] (xreg_split) |- (mul3);
    \draw[->] (w1) -- (mul3);
    \node[below= 0.1cm of mul3,xshift=-0.5cm] (m3) {$M_3$};

    % Adder w2*x^2 + w1*x   
    \node[circop,right=1.3cm of mul2,fill=green!20] (add1) {$+$};
    \draw[->] (mul2) -- node[midway,above] {\texttt{xsq\_w2}} (add1); 
    \draw[->] (mul3) -| node[midway,below] {\texttt{x\_w1}} (add1);

    % Adder + w0
    \node[circop,right=1.5cm of add1,fill=green!20] (add2) {$+$};
    \node [above of=add2, yshift=0.5cm] (w0) {\texttt{w0}};
    \draw[->] (add1) -- (add2);
    \draw[->] (w0) -- (add2);
    \node[below= 0.1cm of add1,xshift=0.5cm] (a1) {$A_1$};
    \node[below= 0.1cm of add2] (a2) {$A_2$};

    % Output
    \draw[->] (add2) -- ++(1,0) node[right] {\texttt{y}};
  
\end{tikzpicture}
  \caption{Single cycle implementation of a quadratic function.}
  \label{fig:quad_func}
\end{figure}

\subsection{Two cycle implementation}
\begin{systemverilogcode}
always_ff @(posedge clk) begin
  x_reg <= x;
  lin_term <= w1 * x_reg + w0;
  xsq <= x_reg * x_reg;
end
always_comb begin
  y = w2 * xsq + lin_term;
end
\end{systemverilogcode}

\subsection{General Pipelining}

Non-pipelined loop body:
\begin{systemverilogcode}
always_ff @(posedge clk) begin
  x_reg <= x;
end
always_comb begin
  z1 = f1(x_reg);
  z2 = f2(z1);
  y = f3(z2);
end
\end{systemverilogcode}

Pipelined loop body:
\begin{systemverilogcode}
always_ff @(posedge clk) begin
  x_s0 <= x;
  z1_s1 <= f1(x_s0);
  z2_s2 <= f2(z1_s1);
end
always_comb begin
  y = f3(z2_s2);
end
\end{systemverilogcode}

Diagram:

\begin{figure}
\centering
\begin{tikzpicture}
    \node (x) {\texttt{x}};
    \node[shape=register, draw, minimum width=0.7cm, minimum height=1cm,
        right=1cm of x,fill=blue!10] (xreg) {};
    \draw[->] (x) -- (xreg.input);
    \node [shape=rectangle, draw, minimum width=1cm, minimum height=1cm,
        right=1cm of xreg,fill=green!20] (f1) {$f_1$};
    \node [shape=rectangle, draw, minimum width=1cm, minimum height=1cm,
      right=1cm of f1,fill=green!20] (f2) {$f_2$};
     \node [shape=rectangle, draw, minimum width=1cm, minimum height=1cm,
      right=1cm of f2,fill=green!20] (f3) {$f_3$};
    
    \draw[->] (xreg) -- (f1);
    \draw[->] (f1) -- node[midway,above] {\texttt{z1}} (f2);
    \draw[->] (f2) -- node[midway,above] {\texttt{z2}} (f3);
    \draw[->] (f3) -- ++(1,0) node[right] {\texttt{y}};
\end{tikzpicture}
  \caption{Unpipelined loop body.}
\end{figure}

\begin{figure}
\centering
\begin{tikzpicture}
    \node (x) {\texttt{x}};
    \node[shape=register, draw, minimum width=0.7cm, minimum height=1cm,
        right=1cm of x,fill=blue!10] (xreg) {};
    \draw[->] (x) -- (xreg.input);
    \node [shape=rectangle, draw, minimum width=1cm, minimum height=1cm,
        right=1cm of xreg,fill=green!20] (f1) {$f_1$};
    \node[shape=register, draw, minimum width=0.7cm, minimum height=1cm,
        right=1cm of f1,fill=blue!10] (f1reg) {};
    \node [shape=rectangle, draw, minimum width=1cm, minimum height=1cm,
      right=1.5cm of f1reg,fill=green!20] (f2) {$f_2$};
    \node[shape=register, draw, minimum width=0.7cm, minimum height=1cm,
        right=1cm of f2,fill=blue!10] (f2reg) {};
    \node [shape=rectangle, draw, minimum width=1cm, minimum height=1cm,
      right=1.5cm of f2reg,fill=green!20] (f3) {$f_3$};
    
    \draw[->] (xreg) -- node[midway,above] {\texttt{x\_s0}} (f1);
    \draw[->] (f1) -- (f1reg);
    \draw[->] (f1reg) -- node[midway,above] {\texttt{z1\_s1}} (f2);
    \draw[->] (f2) -- (f2reg);
    \draw[->] (f2reg) -- node[midway,above] {\texttt{z2\_s2}} (f3);
    \draw[->] (f3) -- ++(1,0) node[right] {\texttt{y}};
\end{tikzpicture}
  \caption{Pipelined loop body.}
\end{figure}

\subsection{Pipelining with multiple paths}


\begin{systemverilogcode}
always_ff @(posedge clk) begin
  x1_reg <= x1;
  x2_reg <= x2;
end
always_comb begin
  z1 = f1(x1_reg);
  z2 = f2(x2_reg);
  z3 = f3(z1, x1_reg);
  y = f4(z2, z3);
end
\end{systemverilogcode}

\pagebreak
A pipeline version of this could be:
\begin{systemverilogcode}
always_ff @(posedge clk) begin
  x1_s0 <= x1;
  x2_s0 <= x2;
  z1_s1 <= f1(x1_s0);
  z2_s1 <= f2(x2_s0);
  x1_s1 <= x1_s0;
  z3_s2 <= f3(z1_s1, x1_s1);
  z2_s2 <= z2_s1;
end
always_comb begin
  y = f4(z2_s2, z3_s2);
end
\end{systemverilogcode}

\begin{figure}
\centering
\begin{tikzpicture}[font=\footnotesize]
    \node (x1) {\texttt{x1}};
    \node[shape=register, draw, minimum width=0.7cm, minimum height=1cm,
        right=1cm of x1,fill=blue!10] (x1reg) {};
    \draw[->] (x1) -- (x1reg.input);

    \node (x2) [below=1cm of x1]{\texttt{x2}};
    \node[shape=register, draw, minimum width=0.7cm, minimum height=1cm,
        right=1cm of x2,fill=blue!10] (x2reg) {};
    \draw[->] (x2) -- (x2reg.input);

    \node [shape=rectangle, draw, minimum width=1cm, minimum height=1cm,
        right=1cm of x1reg,fill=green!20] (f1) {$f_1$};
    \node [shape=rectangle, draw, minimum width=1cm, minimum height=1cm,
      right=1cm of x2reg,fill=green!20] (f2) {$f_2$};
    \node [shape=rectangle, draw, minimum width=1cm, minimum height=1cm,
      right=1cm of f1,fill=green!20] (f3) {$f_3$};
    \node [shape=rectangle, draw, minimum width=1cm, minimum height=1cm,
      right=1cm of f3,fill=green!20] (f4) {$f_4$};
    
    \draw[->] (x1reg) -- ++(0.5,0) -- ++(0,1) -| (f3.north);
    \draw[->] (x1reg) -- (f1);
    \draw[->] (x2reg) -- (f2);
    \draw[->] (f1) -- node[midway,above] {\texttt{z1}} (f3);
    \draw[->] (f2) -| node[midway,above,xshift=-0.5cm] {\texttt{z2}} (f4.south);
    \draw[->] (f3) -- node[midway,above] {\texttt{z3}} (f4);
    \draw[->] (f4) -- ++(1,0) node[right] {\texttt{y}};
\end{tikzpicture}
  \caption{Unpipelined loop body.}
\end{figure}



\begin{figure}
\centering
\begin{tikzpicture}[font=\footnotesize]
    \node (x1) {\texttt{x1}};
    \node[shape=register, draw, minimum width=0.7cm, minimum height=1cm,
        right=1cm of x1,fill=blue!10] (x1reg) {};
    \draw[->] (x1) -- (x1reg.input);

    \node (x2) [below=1cm of x1]{\texttt{x2}};
    \node[shape=register, draw, minimum width=0.7cm, minimum height=1cm,
        right=1cm of x2,fill=blue!10] (x2reg) {};
    \draw[->] (x2) -- (x2reg.input);

    \node [shape=rectangle, draw, minimum width=1cm, minimum height=1cm,
        right=1cm of x1reg,fill=green!20] (f1) {$f_1$};
    \node [shape=rectangle, draw, minimum width=1cm, minimum height=1cm,
      right=1cm of x2reg,fill=green!20] (f2) {$f_2$};
    \node[shape=register, draw, minimum width=0.7cm, minimum height=1cm,
        right=1cm of f1,fill=blue!10] (z1_s1) {};
    \node[shape=register, draw, minimum width=0.7cm, minimum height=1cm,
        right=1cm of f2,fill=blue!10] (z2_s1) {};
    \node[shape=register, draw, minimum width=0.7cm, minimum height=1cm,
        above=0.7cm of z1_s1,fill=blue!10] (x1_s1) {};
    \node [shape=rectangle, draw, minimum width=1cm, minimum height=1cm,
      right=1cm of z1_s1,fill=green!20] (f3) {$f_3$};
    \node[shape=register, draw, minimum width=0.7cm, minimum height=1cm,
        right=1cm of f3,fill=blue!10] (z3_s2) {};
     \node[shape=register, draw, minimum width=0.7cm, minimum height=1cm,
        right=3cm of z2_s1,fill=blue!10] (z2_s2) {};
    \node [shape=rectangle, draw, minimum width=1cm, minimum height=1cm,
      right=1cm of z3_s2,fill=green!20] (f4) {$f_4$};
    
    \draw[->] (x1reg) -- ++(0.7,0) |- node[midway,above,xshift=1cm] {\texttt{x1\_s0}} (x1_s1.input);
    \draw[->] (x1_s1) -| node[midway,above,xshift=-0.5cm] {\texttt{x1\_s1}} (f3.north);
    \draw[->] (x1reg) -- (f1);
    \draw[->] (x2reg) -- (f2);
    \draw[->] (f1) -- (z1_s1); 
    \draw[->] (z1_s1) -- node[midway,above] {\texttt{z1\_s1}} (f3);
    \draw[->] (f2) -- (z2_s1);
    \draw[->] (z2_s1) -- node[midway,above] {\texttt{z2\_s1}} (z2_s2.input);
    \draw[->] (z2_s2) -| node[midway,above,xshift=-0.5cm] {\texttt{z2\_s2}} (f4.south);
    \draw[->] (f3) --  (z3_s2);
    \draw[->] (z3_s2) -- node[midway,above] {\texttt{z3\_s2}} (f4);
    \draw[->] (f4) -- ++(1,0) node[right] {\texttt{y}};
\end{tikzpicture}
  \caption{Multi-path with pipelining.}
\end{figure}

\pagebreak
\subsection*{Memory Accesses}

Consider the following loop to compute the product of two vectors:
\begin{systemverilogcode}
logic [31:0] x[0:N-1], w[0:N-1];
logic [31:0] y[0:N-1];

logic [8:0] i;
always_ff @(posedge clk) begin
  // Some start-stop logic
  // running = ...

  // Stage 0:  Read x[i] and w[i]
  if ((running) && (i < N)) begin
    x_s0 <= x[i];
    w_s0 <= w[i];
  end

  // Stage 1:  Compute product
  if ((running) && (i >= 1) && (i < N+1)) begin
    prod_s1 <= x_s0 * w_s0;
  end

  // Stage 2:  Write output
  if ((running) && (i >= 2) && (i < N+2)) begin
    y[i-2] <= prod_s1;
  end

  if ((running) && (i < N+2)) begin
    i <= i + 1;
  end 
end
\end{systemverilogcode}

The C version would be:

\begin{ccode}
for (int i=0; i < N; i++) {
#pragma HLS PIPELINE II=1
  y[i] = w[i]*x[i];
}
\end{ccode}

\subsection{Polynomial example}

we want to build the equivalent of:
\begin{ccode}
for (i=0; i < n; i++) {
  y[i] = a[0] + a[1]*x[i] 
      + a[2]*x[i]*x[i] + a[3]*x[i]*x[i]*x[i];
}
\end{ccode}

This is most efficiently performed by Horner's method:
\begin{ccode}
for (i=0; i < n; i++) {
  xi = x[i];
  yi = a[3];
  yi += xi*yi + a[2];
  yi += xi*yi + a[1];
  yi += xi*yi + a[0];
  y[i] = yi;
}
\end{ccode}

\subsection*{Loop carry dependency}

We can illustrate loop carry dependency with a merge sort:
First consider non-pipelined version
\begin{systemverilogcode}
always_ff @(posedge clk) begin
  if (start) begin
    i <= 0;
    j <= 0;
    k <= 0;
  end else if (i < N1 && j < N2) begin
    if (x1[i] < x2[j]) begin
      y[k] <= x1[i];
      i <= i + 1;
    end else begin
      y[k] <= x2[j];
      j <= j + 1;
    end
    k <= k + 1;
  end else if (i < N1) begin
    y[k] <= x1[i];
    i <= i + 1;
    k <= k + 1;
  end else if (j < N2) begin
    y[k] <= x2[j];
    j <= j + 1;
    k <= k + 1;
  end
 
end

\end{systemverilogcode}

\begin{ccode}
for (i=0, j=0, k=0; i < N1 && j < N2; k++) {
  if (i == N1) {
    y[k] = x2[j];
    j++;
  } else if (j == N2) {
    y[k] = x1[i];
    i++;
  } else if (x1[i] < x2[j]) {
    y[k] = x1[i];
    i++;
  } else {
    y[k] = x2[j];
    j++;
  }
}
\end{ccode}

Focusing on the main case, 
we need the following steps, each taking a clock cycle
\begin{systemverilogcode}
// stage 0
x1i <= x1[i];
x2j <= x2[j];

// stage 1
if (x1i < x2j) begin
  yk <= x1i;
  i <= i + 1;
end else begin
  yk <= x2j;
  j <= j + 1;
end

// stage 2
y[k] <= yk;
k <= k+1;

\end{systemverilogcode}

Stage 0 of iteration $t$ depends on the results of stage 1 of iteration $t-1$.


\subsection*{Loop unrolling}

\begin{ccode}
// Normal loop
for (i=0; i < N; i++) {
  y[i] = w[i]*x[i];
}

// Unrolled by factor of 2
for (i=0; i < N; i+=2) {
  y[i] = w[i]*x[i];
  y[i+1] = w[i+1]*x[i+1];
}

\end{ccode}

\begin{ccode}
// Vitis pragma
for (i=0; i < N; i++) {
#pragma HLS unroll factor=2
  y[i] = w[i]*x[i];
}

\end{ccode}

But this runs into a memory bottleneck.

\begin{systemverilogcode}

// Stage 0:  Two reads from x and w
x1_s0 <= x[i];
x2_s0 <= x[i+1];
w1_s0 <= w[i];
w2_s0 <= w[i+1];

// Stage 1:  Two products  -> OK
y1_s1 <= x1_s0 * w1_s0;
y2_s1 <= x2_s0 * w2_s0;

// Stage 2:  Two writes
y[i] <= y1_s1;
y[i+1] <= y2_s1;
i <= i + 2;
\end{systemverilogcode}

\subsection*{Memory bottleneck}

\begin{ccode}

ap_uint<64> xwords[N/2], wwords[N/2], ywords[N/2];

for (i=0; i < N; i += 2) {

  // Read two words from x and w
  xword = xwords[i/2];
  wword = wwords[i/2];

  // Exact values
  x1 = xword.range(31,0);
  x2 = xword.range(63,32);
  w1 = wword.range(31,0);
  w2 = wword.range(63,32);

  // Compute products
  y1 = w1*x1;
  y2 = w2*x2;

  // Write outputs
  yword.range(31,0) = y1;
  yword.range(63,32) = y2;
  ywords[i/2] = yword;
}
\end{ccode}
\end{document}

  

